\documentclass[a4paper,11pt]{article}
\usepackage{amssymb, enumitem}
\setlist[enumerate,1]{label={\bfseries\arabic*.},ref={\bfseries{Ejercicio \arabic*}}}
\setlist[enumerate,2]{label={\bfseries(\alph*)},ref={\bfseries(\alph*)}}


\parindent 0cm
\usepackage{amssymb,amsmath,amsthm,latexsym,epsfig,euscript,multicol}
\usepackage{graphbox}
\usepackage{booktabs}

\usepackage[utf8x]{inputenc}
\usepackage{listings,xcolor,bm}
\usepackage{url}
\usepackage{hyperref}
\hypersetup{
    colorlinks=true,
    urlcolor=blue,
    }


\definecolor{mygreen}{rgb}{0,0.6,0}
\definecolor{mygray}{rgb}{0.5,0.5,0.5}
\definecolor{mymauve}{rgb}{0.58,0,0.82}
\lstset{
  backgroundcolor=\color{white},   % choose the background color; you must add
  basicstyle=\ttfamily,      % the size of the fonts that are used for the code
  breakatwhitespace=true,         % sets if automatic breaks should only happen at whitespace
  breaklines=true,                 % sets automatic line breaking
  captionpos=b,                    % sets the caption-position to bottom
  commentstyle=\color{mygreen},    % comment style
  deletekeywords={...},            % if you want to delete keywords from the given language
  escapeinside={\%*}{*)},          % if you want to add LaTeX within your code
  extendedchars=true,              % lets you use non-ASCII characters; for 8-bits encodings only, does not work with UTF-8
  firstnumber=1,                % start line enumeration with line 1000
  frame=none,	                   % adds a frame around the code
  keepspaces=true,                 % keeps spaces in text, useful for keeping indentation of code (possibly needs columns=flexible)
  keywordstyle=\color{blue},       % keyword style
  language=Python,                 % the language of the code
  morekeywords={*,...},            % if you want to add more keywords to the set
  numbers=none,                    % where to put the line-numbers; possible values are (none, left, right)
  numbersep=5pt,                   % how far the line-numbers are from the code
  numberstyle=\tiny\color{mygray}, % the style that is used for the line-numbers
  rulecolor=\color{black},         % if not set, the frame-color may be changed on line-breaks within not-black text (e.g. comments (green here))
  showspaces=false,                % show spaces everywhere adding particular underscores; it overrides 'showstringspaces'
  showstringspaces=false,          % underline spaces within strings only
  showtabs=false,                  % show tabs within strings adding particular underscores
  stepnumber=5,                    % the step between two line-numbers. If it's 1, each line will be numbered
  stringstyle=\color{mymauve},     % string literal style
  tabsize=4,	                   % sets default tabsize to 2 spaces
  title=\lstname,                  % show the filename of files included with \lstinputlisting; also try caption instead of title
  belowskip=0em
}
% Caracteres especiales
\def\A{\mathbb{A}}
\def\C{\mathbb{C}}
\def \N{\mathbb{N}}
\def \P{\mathbb{P}}
\def \Q{\mathbb{Q}}
\def \R{\mathbb{R}}
\def \Z{\mathbb{Z}}
\def \sen{\textrm{sen}}

\def\Np{$\N$}
\def\Zp{$\Z$}
\def\Qp{$\Q$}
\def\Rp{$\R$}
\def\Cp{$\C$}

\def\bb{\bm{b}}
\def\bu{\bm{u}}
\def\bv{\bm{v}}
\def\bx{\bm{x}}
\def\bA{\bm{A}}
\def\bB{\bm{B}}
\def\bD{\bm{D}}
\def\bE{\bm{E}}
\def\bM{\bm{M}}
\def\bT{\bm{T}}


\def\K{\textrm{K}}
\def\V{\textrm{V}}
\def\S{\textrm{S}}

\def\degres{$^\circ$}

\newcount\todno
\def\no{\global\advance\todno by 1 \the\todno}

\topmargin-2cm \vsize 29.5cm \hsize 21cm
\setlength{\textwidth}{16.75cm}\setlength{\textheight}{23.5cm}
\setlength{\oddsidemargin}{0.0cm}
\setlength{\evensidemargin}{0.0cm}


\theoremstyle{definition}
\newtheorem{ejer}{Ejercicio}
\newcommand{\bej}{\begin{ejer}}
\newcommand{\fej}{\end{ejer}}

\begin{document}

\centerline{{\small Universidad de Buenos Aires - Facultad de Ciencias Exactas y Naturales - Ciencias de Datos}}

\vskip 0.2cm

\hrule

\vskip 0.2cm

 \centerline{{\bf\Large{\sc Laboratorio de Datos}}}

 \vskip 0.2cm

 \centerline{\ttfamily Primer Cuatrimestre 2024}

\vskip 0.2cm

 \hrule

 \bigskip
 \centerline{\bf Trabajo Práctico N$^\circ$ 2}
 \bigskip

El Trabajo Práctico deberá ser resuelto en grupos de dos o tres personas. No se aceptarán entregas individuales. La entrega se realizará a través del campus (pestaña Trabajos Prácticos). La fecha límite es el 28/05 a las 23:59. Deben entregar un Notebook con los nombres de les integrantes del equipo, la resolución de los ejercicios y los informes pertinentes.

Se valorará que el Notebook y el código tengan un formato prolijo: ejercicios separados por títulos (Ejercicio 1, Ejercicio 2, etc.), nombres descriptivos para las variables, comentarios, etc.

\vspace{0.5cm}

En este trabajo práctico nos sumamos a la moda de utilizar an\'alisis de datos para la toma de decisiones en los deportes, y creamos nuestra propia empresa LDD Futbol Analytics.


Trabajaremos con el dataset \verb|FBRef2020-21.csv|\footnote{Fuente: \url{https://fbref.com/en/}} que contiene datos sobre los principales torneos de clubes y selecciones de fútbol del mundo.


\subsection*{Preprocesamiento [2 pt.]}

\subsection*{Clustering [4 pts.]}

\begin{enumerate}[resume]
\item Nuestro primer objetivo es realizar algún agrupamiento de jugadores con características similares.
\begin{enumerate}
\item Seleccionar dos variables cualesquiera de los datos y realizar un gráfico de dispersión de una variable en función de la otra para el total de las observaciones. ¿Pueden encontrar fácilmente grupos distintos?

\item Escalar los datos y realizar un análisis de componentes principales, quedándose solo con las dos primeras componentes.
Realizar un gráfico como el del punto anterior. ¿Cuántos clusters puede distinguir en el gráfico? ¿A qué características de los jugadores pueden corresponder los clusters? ¿Cómo pueden verificar su conjetura? (realizar una visualización o algún cálculo)

\item Para la cantidad de clusters observados en el \'item anterior, realizar un agrupamiento por $k$-medias, y colorear los puntos según las etiquetas obtenidas. Coinciden las etiquetas con lo esperado?

\item Repetir el agrupamiento utilizando DBSCAN. ¿Cómo eligirían en este caso un valor de $\varepsilon$ apropiado? Sugerencia: consultar la sección "Selección del hiperparámetro eps" del Notebook de la clase de DBSCAN (o utilizar cualquier otra técnica que consideren apropiada)

\item Para realizar el clustering pueden utilizar como datos todas las variables originales o solo las dos componentes principales. Si en el item anterior utilizaron una de estas alternativas, repitan el experimento utilizando la otra alternativa (modificando los valore de epsilon y minPts convenientemente). ¿Con cuál de las dos opciones obtienen mejores resultados?

\end{enumerate}
\end{enumerate}

\subsection*{Clasificación [3 pts.]}

\begin{enumerate}[resume]
\item Ahora queremos poder predecir la posición en la que juega cada jugador según sus datos estadísticos utilizando $KNN$. En la columna Pos encontramos la posición de los jugadores. Para la mayoría de los jugadores se indica una \'unica posición pero algunos jugadores tienen dos posiciones. Para simplificar el análisis vamos a considerar una única posición por jugador.
\begin{enumerate}
\item Elegir alguna opción para quedarse con una única posición por jugador (pueden quedarse solo con la primera o eliminar a los jugadores con dos posiciones o alguna otra estrategia que se les ocurra).
\item Construir una base de datos $X$ con las variables explicativas e $y$ con las variable de posición.
\item Dividir ambos conjunto de datos en un 80\% para entrenamiento y un 20\% para testeo.
\item Aplicar un esquema de validación en el conjunto de entrenamiento para seleccionar el valor óptimo de $K$. (Esto puede demorar mucho si prueban mucho, hacerlo hasta un valor máximo de $K = 20$.)
\item Para el valor de $K$ obtenido, ¿cuál es el porcentaje de aciertos en el conjunto de testeo?
\item Al igual que en el caso de clustering, pueden realizar la clasificación utilizando todas las variables o solo las primeras componentes principales. Utilizar ambos métodos e indicar con cuál método obtienen mejores resultados.
\end{enumerate}
\end{enumerate}

\subsection*{Redes neuronales [1 pts.]}

\begin{enumerate}[resume]
\item Trabajamos ahora con el dataset de jugadoras de la 2023-2024 Women's Super League de Inglaterra.
¿Cómo podemos predecir la posición en qué juega cada jugadora utilizando redes neuronales? ¿Qué porcentaje de aciertos obtienen?
\end{enumerate}


\end{document}

